\documentclass[10pt]{article}
\usepackage[T1]{fontenc}
\usepackage[utf8]{inputenc}
\usepackage[french]{babel}
\usepackage{amsmath,amssymb,mathrsfs,mathpazo}
\usepackage{geometry}
\geometry{a4paper,margin=2cm}
\usepackage{graphicx,array,tabularx,tabularray,colortbl,multirow,multicol,booktabs,enumitem,pifont,ulem,fancybox,diagbox,cancel}
\usepackage{pstricks,pstricks-add,pst-plot,pst-tree,pst-node,pst-text,pst-func,pst-math,pst-bspline,pst-3dplot,pst-eucl,pst-solides3d}
\let\pscalculate\undefined
\usepackage{pst-bezier}
\usepackage{tikz,pgfplots,tkz-tab}
\usetikzlibrary{arrows,arrows.meta,patterns,calc,positioning,trees,decorations.markings,shapes.geometric,shapes.misc,shapes.arrows,intersections,decorations.pathmorphing,decorations.pathreplacing,backgrounds,shadows,babel}
\pgfplotsset{compat=1.18}
\usepgfplotslibrary{fillbetween}
\usepackage[np]{numprint}
\usepackage{listings,makecell,pas-tableur,twoopt}
\let\barre\undefined
\usepackage[pstricks]{tablvar}
\definecolor{darkgreen}{RGB}{0,100,0}
\providecommand{\red}{\color{red}}
\providecommand{\blue}{\color{blue}}
\providecommand{\green}{\color{green}}
\providecommand{\black}{\color{black}}
\providecommand{\white}{\color{white}}
\providecommand{\gray}{\color{gray}}
\providecommand{\cyan}{\color{cyan}}
\providecommand{\magenta}{\color{magenta}}
\providecommand{\yellow}{\color{yellow}}
\providecommand{\violet}{\color{violet}}
\providecommand{\orange}{\color{orange}}
\providecommand{\euro}{\texteuro}
\providecommand{\EUR}{\texteuro}
\providecommand{\ieme}{\textsuperscript{e}}
\providecommand{\degres}{\ensuremath{^\circ}}
\providecommand{\degree}{\ensuremath{^\circ}}
\providecommand{\vv}[1]{\overrightarrow{#1}}
\providecommand{\Coord}[1]{#1}
\providecommand{\decosix}{}
\providecommand{\starredbullet}{\ensuremath{\bullet}}
\pagestyle{empty}
\setlength{\parindent}{0pt}
\usepackage{adjustbox}\DeclareUnicodeCharacter{2236}{:}\DeclareUnicodeCharacter{2190}{\ensuremath{\leftarrow}}
\providecommand{\R}{\mathbb{R}}
\providecommand{\N}{\mathbb{N}}
\providecommand{\Z}{\mathbb{Z}}
\providecommand{\Q}{\mathbb{Q}}
\providecommand{\D}{\mathbb{D}}
\providecommand{\C}{\mathbb{C}}
\providecommand{\e}{{\mathrm{e}}}
\providecommand{\vect}[1]{\overrightarrow{#1}}
\providecommand{\vectt}[1]{\overrightarrow{\mathrm{#1}}}
\providecommand{\vvt}[1]{\overrightarrow{\mathrm{#1}}}
\providecommand{\np}[2][]{{#2~#1}}
\providecommand{\vv}[1]{{\overrightarrow{#1}}}
\renewcommand{\d}{{\mathrm{d}}}
\providecommand{\degres}{{^\circ}}
\providecommand{\degre}{{^\circ}}
\providecommand{\blue}{\color{blue}}
\providecommand{\red}{\color{red}}
\providecommand{\enskip}{\quad}
\providecommand{\llbracket}{\lbrack\!\lbrack}
\providecommand{\rrbracket}{\rbrack\!\rbrack}
\providecommand{\euro}{{\text{€}}}
\providecommand{\no}{{\text{n°}}}
\providecommand{\interval}[2]{{[#1;#2]}}
\providecommand{\intervFO}[2]{{[#1;#2[}}
\providecommand{\dotfill}{\ldots}

\providecommand{\ds}{\displaystyle}
\providecommand{\textding}[1]{#1}
\providecommand{\Oij}{\(\mathrm{O};\vec{\imath},\vec{\jmath}\)}
\providecommand{\Oijk}{\(\mathrm{O};\vec{\imath},\vec{\jmath},\vec{k}\)}
\providecommand{\Ouv}{\(\mathrm{O};\vec{u},\vec{v}\)}
\providecommand{\intervFF}[2]{\ensuremath{[#1;#2]}}
\renewcommand{\euro}{\ifmmode\text{\texteuro}\else\texteuro\fi}
\DeclareUnicodeCharacter{00B0}{\ifmmode^{\circ}\else\textdegree\fi}
\DeclareUnicodeCharacter{20AC}{\ifmmode\text{\texteuro}\else\texteuro\fi}
% Énoncés : Éducation nationale. Transcription et corrigés : APMEP, Denis Vergès et contributeurs. Découpage et classement : CoopMaths, Sébastien Lozano, Éric Elter et contributeurs. Adaptation de la présentation : Maths Explorer.
% CC BY-SA : https://coopmaths.fr/www/about/
% https://coopmaths.fr/alea/static/bac/2022/tex/bac_2022_05_sujet1_metropole_2_cor.tex
\begin{document}
\noindent\textit{APMEP / CoopMaths. CC BY-SA. Présentation adaptée par Maths Explorer.}\par\medskip
\medskip

\begin{enumerate}
	\item 
	\begin{enumerate}
	\item La droite $\mathcal{D}$ dont une représentation paramétrique est $\left\{\begin{array}{l c l}
		x&=&1+2t\\y &=& 2 - t\\z&=& 2+2t
	\end{array}\right.  , t \in \R$ a pour vecteur directeur \(\overrightarrow{u}\begin{pmatrix}2\\-1\\2\end{pmatrix}\).
	\item Si $B$ appartient à la droite $\mathcal{D}$ alors $\exists t \in \R$ tel que :
	
	 $\left\{\begin{array}{l c l} -1&=&1+2t\\3 &=& 2 - t\\0&=& 2+2t
	\end{array}\right.\iff \left\{\begin{array}{l c l}
	2t&=&-2\\-t&=&1\\2t&=&-2t\end{array}\right. \iff t=-1 $. Donc $B \in  \mathcal{D}$.
\item On a $A(-1~;~1~;~3)$ et $B(-1~;~3~;~0)$. Donc \(\overrightarrow{AB}\begin{pmatrix}0\\2\\-3\end{pmatrix}\).

Alors $\vect{AB}.\vect{u}=0\times 2+2\times (-1)+(-3)\times 2=-8$
\end{enumerate}
\item 
	\begin{enumerate}
		\item Le plan $\mathcal{P}$ passant par $A$ est orthogonal à la droite $\mathcal{D}$. Donc $\mathcal{P}$ a pour vecteur normal le vecteur $\vect{u}$ (vecteur directeur de $\mathcal{D}$).
	
Son équation cartésienne sera de la forme $ax+by+cz+d=0$, où $(a~;~b~;~c)$ sont les cordonnées d'un vecteur normal au plan.
	
En prenant comme vecteur normal à $\mathcal{P}$ le vecteur $\vect{u}$, on obtient : $\mathcal{P}~:$

$2x-y+2z+d=0$. Or $A\in \mathcal{P}$ donc $-2-1+6+d=0$ donc $d=-3$. Donc $\mathcal{P}$ a pour équation $2x-y+2z-3=0$
		\item Le projeté orthogonal de $A$ sur la droite $\mathcal{D}$, noté $H$, est l'unique point d'intersection de $\mathcal{P}$ et $\mathcal{D}$. Résolvons donc le système suivant : \\
	$\left\{\begin{array}{l c l} 2x-y + 2z-3&=&0\\x&=&1+2t\\y &=& 2 - t\\z&=& 2+2t \end{array}\right.$. On remplace $x$, $y$ et $z$ dans la première équation :\\
	$2x-y+2z-3=0\iff 2(1 + 2t)-(2 - t)+2(2 + 2t) - 3 = 0 \iff $
	
$2 + 4t  - 2 + t + 4 + 4t = 0 \iff 9t + 1= 0\iff t = -\dfrac{1}{9}$.\\
	On remplace la valeur de $t$ dans les trois dernières équations :\\
		$\left\{\begin{array}{l c l} x&=&1+2\times-\dfrac{1}{9}=\dfrac{7}{9} \\y &=& 2 - \left(-\dfrac{1}{9}\right)=\dfrac{19}{9}\\z&=& 2+2 \times \left(-\dfrac{1}{9}\right) =\dfrac{16}{9} \end{array}\right.$. Donc $H\left(\dfrac79~;~\dfrac{19}{9}~;~\dfrac{16}{9}\right)$.
		\item Déterminons les coordonnées de $\vect{AH}$ avec $A(-1~;~1~;~3)$ et $H\left(\dfrac79~;~\dfrac{19}{9}~;~\dfrac{16}{9}\right)$ : \(\overrightarrow{AH}\begin{pmatrix}\frac{16}{9}\\\frac{10}{9}\\-\frac{11}{9}\end{pmatrix}\).

Donc $\left\| \vect{AH}\right\| =\sqrt{\left(\dfrac{16}{9}\right)^2+\left(\dfrac{10}{9}\right)^2+\left(-\dfrac{11}{9}\right)^2}=\dfrac{\sqrt{477}}{9}=\dfrac{\sqrt{53}}{3}$
\end{enumerate}
\item 
\begin{enumerate}
	\item Les points $H$ et $B$ appartiennent à $\mathcal{D}$, donc le vecteur $\vect{HB}$ est un vecteur directeur de $\mathcal{D}$, tout comme $\vect{u}$. Donc les vecteurs $\vect{HB}$ et $\vect{u}$ sont colinéaires donc $\exists k \in \R$ tel que \\$\vect{HB}=k\times \vect{u}$.
	\item D'après les propriétés du produit scalaire, et en utilisant la relation de Chasles, \\ $\vect{AB}.\vect{u}=\left(\vect{AH}+\vect{HB}\right).\vect{u}=\vect{AH}.\vect{u}+ \vect{HB}.\vect{u}$.
	
Or les points $A$ et $H$ appartiennent au plan $\mathcal{P}$ normal à la droite $\mathcal{D}$, donc tout vecteur de $\mathcal{P}$ est orthogonal à tout vecteur de $\mathcal{D}$ donc $\vect{AH}\perp \vect{u}$ donc $\vect{AH}.\vect{u}=0$.

Donc $\vect{AB}.\vect{u}=\vect{HB}.\vect{u} = k \times \vect{u}.\vect{u}=k~\vect{u}^2=k~\lVert\vect{u}\rVert^2$ donc $k=\dfrac{\vect{AB}.\vect{u}}{\lVert\vect{u}\rVert^2}$.
	\item On a \(\overrightarrow{u}\begin{pmatrix}2\\-1\\2\end{pmatrix}\) donc $\lVert\vect{u}\rVert=\sqrt{2^2+(-1)^2+2^2}=\sqrt{9}=3$ et $\vect{AB}.\vect{u}=-8$ donc 
	
$k=\dfrac{-8}{3^2}=-\dfrac{8}{9}$.

Donc en posant $H(x~;~y~;~z)$ et $B(-1~;~3~;~0)$ alors \(\overrightarrow{HB}\begin{pmatrix}- 1 - x\\3 - y\\- z\end{pmatrix}\).

De plus $\vect{HB}=k\times \vect{u}=-\dfrac{8}{9}\times \vect{u}$ donc 	$\left\{\begin{array}{l c l}- 1 - x&=&-\dfrac{8}{9}\times 2\\3 - y &=& -\dfrac{8}{9}\times (- 1)\\-z&=&-\dfrac{8}{9}\times 2 \end{array}\right.$

Donc $x=\dfrac{16}{9}-1=\dfrac{7}{9}$ , $y=3-\dfrac{8}{9}=\dfrac{19}{9}$ et $z=\dfrac{16}{9}$. On retrouve les coordonnées du point $H$.
\end{enumerate}
		\item Les points $A$, $H$ et $C$ appartiennent au plan $\mathcal{P}$. $H$ est le projeté orthogonal de B sur  $\mathcal{P}$. Le tétraèdre $BAHC$ a pour base le triangle $AHC$ et pour hauteur $BH$.

Donc $\mathcal{V}_{BAHC}=\dfrac{1}{3}\times \mathcal{A}_{AHC}\times BH$ d'où $\mathcal{A}_{AHC}=\dfrac{3\times \mathcal{V}_{BAHC}}{BH}$.

D'après la question \textbf{3. c}, $\vect{HB}=k\times \vect{u}=$ donc \(\overrightarrow{HB}\begin{pmatrix}\frac{16}{9}\\-\frac{8}{9}\\\frac{16}{9}\end{pmatrix}\) donc

$HB=\sqrt{\left(\dfrac{16}{9}\right)^2+\left(-\dfrac{8}{9}\right)^2+\left(\dfrac{16}{9}\right)^2} = \dfrac{8}{3}$

Donc  $\mathcal{A}_{AHC}=\dfrac{3\times \dfrac{8}{9}}{\dfrac{8}{3}}=1$ unité d'aire.
\end{enumerate}

\bigskip

\end{document}
